\section{\bi{中央门槛下矩形几何的限制}{A geometric obstruction at central thresholds}}
\bi{上界模型允许后一个动作随前一步变化；要改进它，须利用同一接收点的一次成功由同一个矩形实现。一个可显式核算的情形是两个坐标块上的嵌套乘积族。在固定矩形上独立选两块中的均匀点，每一块上的嵌套事件可由同一均匀坐标的阈值表示。下面确定由有限多个角点逼近一个高斯半平面所损失的概率阶。}{The upper model allows an action to depend on preceding increments. Further geometric information comes from the requirement that a receiver be witnessed by a single rectangle. For nested product families on two coordinate blocks, independent uniform sampling in the factors of a fixed rectangle represents each nested family by thresholds on one uniform coordinate. The following computation determines the order of the probability lost when finitely many corners approximate a Gaussian half-plane.}
\begin{lem}
\bi{对独立标准正态 $X,Y$，令 $s,t>0$，$s^2+t^2=1$。取 $z$ 于固定紧区间。以 $M$ 个包含于 $\{sX+tY\ge z\}$ 的上象限之并从内部逼近该半平面。最小概率损失介于 $cst/M$ 与 $1/(M+1)$ 之间。}{Let $X,Y$ be independent standard normals, $s,t>0$, $s^2+t^2=1$, and let $z$ range over a fixed compact interval. The infimum of the missing probability when the half-plane $\{sX+tY\ge z\}$ is approximated from within by $M$ upper quadrants lies between $cst/M$ and $1/(M+1)$.}
\end{lem}
\begin{proof}
\bi{作正交变换 $T=sX+tY,\ S=tX-sY$。在 $T=z+h$ 的切线上，每一象限沿 $S$ 的截长至多 $h/(st)$：由角点沿两轴的非负位移满足加权和至多 $h$。在 $0<h<st/M$ 内，$[-1,1]$ 中至少有长度 1 未覆盖。两个旋转坐标独立，在该紧带内密度有统一正下界，积分给出下界。}{Rotate to $T=sX+tY,\ S=tX-sY$. On $T=z+h$ the section of one quadrant has length at most $h/(st)$ in $S$, since the nonnegative displacements from its corner have weighted sum at most $h$. For $0<h<st/M$ at least unit length in $[-1,1]$ remains uncovered. The joint density in this compact strip is bounded below, giving the lower bound by integration.}

\bi{上界改用均匀坐标 $u=\Phi(X),v=\Phi(Y)$。半平面的下边界为递减函数 $f(u)=\Phi((z-s\Phi^{-1}(u))/t)$，端值为 1 与 0。取 $u_j=j/(M+1)$，$1\le j\le M$，角点为 $(u_j,f(u_j))$。加上虚拟端点 $u_0=0,u_{M+1}=1$，第 $j$ 条竖带内漏失面积至多为 $[f(u_j)-f(u_{j+1})]/(M+1)$。求和即得。}{For the upper bound pass to uniform coordinates $u=\Phi(X),v=\Phi(Y)$. The lower boundary is the decreasing function $f(u)=\Phi((z-s\Phi^{-1}(u))/t)$ with endpoint values one and zero. Use corners $(u_j,f(u_j))$, $u_j=j/(M+1)$, $1\le j\le M$. With virtual endpoints $u_0=0,u_{M+1}=1$, the missing area in strip $j$ is at most $[f(u_j)-f(u_{j+1})]/(M+1)$. The sum telescopes.}
\end{proof}
\bi{另一方面，在单位正方形内对任意多个向下矩形的各边际向下舍入至多 $\Delta$，损失面积至多 $2\Delta$：将原并集整体向左、向下各平移 $\Delta$，留在正方形的部分包含于舍入后的并集。于是平衡块上的 $1/M$ 损失来自半平面的几何逼近。它说明为何固定复杂度的中央估计不能直接统一推广到任意多分量源。}{In contrast, rounding both marginals of any number of lower rectangles down by at most $\Delta$ loses at most $2\Delta$ in area. Translate their original union left and down by $\Delta$; the part remaining in the unit square lies inside the rounded union. Thus the order $1/M$ for balanced blocks comes from the half-plane approximation. Central estimates at fixed complexity consequently require additional information to extend to unrestricted unions.}
